\documentclass[11pt]{article}
\usepackage[margin=1in]{geometry}
\usepackage{amsmath,amssymb}
\usepackage[colorlinks=true,urlcolor=blue,citecolor=blue]{hyperref}
\title{A counterexample to Fisk's totally stable matrix conjecture}
\author{Principia Math in Progress}
\date{12 September 2026}
\begin{document}
\maketitle

\paragraph{Task 1. Resolve the totally stable matrix conjecture.}
Fisk's Conjecture 21.84 in \cite{book}, printed page 648, asserts that
if \(f(x,y)=\sum_{j\ge0}f_j(x)y^j\) is a real polynomial with positive
coefficients, nonvanishing when both variables lie in the open upper
half-plane, then every minor of
\[
 T_f(x)=
 \begin{pmatrix}
 f_0&f_1&f_2&f_3&\cdots\\
 0&f_0&f_1&f_2&\cdots\\
 0&0&f_0&f_1&\cdots\\
 \vdots&\vdots&\vdots&\ddots&\ddots
 \end{pmatrix}
\]
is stable. The same assertion appears as Question 10 of \cite{questions}.
Here a univariate polynomial is \emph{stable} if it has no zero in the
open right half-plane. The conjecture is false.

\paragraph{An admissible input.}
Take
\begin{align*}
 f(x,y)&=(10x+y+1)(x+y+10)\\
       &=10x^2+11xy+y^2+101x+11y+10.
\end{align*}
Every displayed coefficient is positive. If
\(\operatorname{Im}x,\operatorname{Im}y>0\), the imaginary parts of the
two factors are \(10\operatorname{Im}x+\operatorname{Im}y>0\) and
\(\operatorname{Im}x+\operatorname{Im}y>0\), respectively.
Neither factor vanishes. Thus \(f\) belongs to the source's class
\(\mathbf P^{\mathrm{pos}}_2\).
It is also nonvanishing when both variables lie in the open right
half-plane, directly from the real parts of its factors.

\paragraph{The offending minor.}
Write
\[
 a(x)=11x+11,\qquad b(x)=10x^2+101x+10.
\]
Using zero-based indices, choose rows \(0,1,2,3\) and columns \(1,2,3,4\)
of \(T_f\). This gives the genuine minor
\[
 M(x)=
 \begin{pmatrix}
 a&1&0&0\\
 b&a&1&0\\
 0&b&a&1\\
 0&0&b&a
 \end{pmatrix}.
\]
Expansion of these tridiagonal determinants gives \(D_0=1\), \(D_1=a\)
and \(D_k=aD_{k-1}-bD_{k-2}\). Consequently,
\begin{align}
 D(x):=\det M(x)
 &=a^4-3a^2b+b^2 \notag\\
 &=11111x^4+16661x^3+17661x^2+16661x+11111. \label{eq:D}
\end{align}

\paragraph{An exact certificate of failure.}
Define
\[
 G(t)=11111t^2+16661t-4561.
\]
For \(x\ne0\), equation \eqref{eq:D} gives the identity
\begin{equation}\label{eq:reciprocal}
 D(x)=x^2G(x+x^{-1}).
\end{equation}
The exact evaluations
\[
 G(1/5)=-\frac{19609}{25}<0,\qquad
 G(1/4)=\frac{4779}{16}>0
\]
show that \(G\) has a real zero \(t_0\in(1/5,1/4)\).
Set
\[
 z_0=\frac{t_0+i\sqrt{4-t_0^2}}2.
\]
Then \(|z_0|=1\), \(z_0+z_0^{-1}=t_0\), and
\[
 \frac1{10}<\operatorname{Re}z_0<\frac18.
\]
Equation \eqref{eq:reciprocal} proves \(D(z_0)=0\).
Thus this minor has a zero strictly in the forbidden half-plane.
Its conjugate is another such zero. The input satisfies every hypothesis,
so the universal assertion is disproved. This proof uses exact integer
identities and the intermediate value theorem; no approximate root
calculation or boundary convention is needed.

\paragraph{Task 2. Determine the smallest failing minor order.}
Decide whether every minor of order at most three is stable for every
\(f\in\mathbf P^{\mathrm{pos}}_2\), or construct a smaller counterexample.
The example above gives an upper bound of four on the minimum order.

\paragraph{Task 3. Classify the two-factor case.}
For \(f(x,y)=(a_1x+b_1y+c_1)(a_2x+b_2y+c_2)\) with positive parameters,
characterize when every minor of \(T_f(x)\) is stable. Identify how the
first failing order depends on the parameter ratios.

\paragraph{Task 4. Test the weaker coefficient assertion.}
Determine whether every minor of \(T_f(x)\) has nonnegative coefficients
for arbitrary \(f\in\mathbf P^{\mathrm{pos}}_2\). The present counterexample
retains coefficient positivity, so it does not decide that weaker claim.

\paragraph{Task 5. Find valid preservation hypotheses.}
Find structural conditions on polynomial matrices under which
multiplication preserves stability of every minor, and use them to
identify subclasses for which Fisk's original conclusion does hold.

\begin{thebibliography}{9}
\bibitem{book}
S.~Fisk, \emph{Polynomials, roots, and interlacing}, version 2,
\href{https://arxiv.org/abs/math/0612833}{arXiv:math/0612833},
2008, \S21.1--21.2 and Conjecture 21.84.
\bibitem{questions}
S.~Fisk, \emph{Questions about determinants and polynomials},
\href{https://arxiv.org/abs/0808.1850}{arXiv:0808.1850}, 2008,
Definition 1 and Question 10.
\end{thebibliography}
\end{document}
